% ==========================================================================
% GENERATED by scripts/split_paper.py. Do not edit this file.
%
% Source: imported_papers/sds_hounsfield_features/main.tex
% Part 8: Comparing to HU Windowing
%
% Edits belong upstream, in the published paper. To change what the
% chapter includes, change imports/hounsfield.toml and re-run
% the splitter.
% ==========================================================================
\section{Comparing to HU Windowing}\label{sec:hu-windowing}

HU windowing is an alternative pre-processing technique proposed by \cite{kloenne_domain-specific_2020}, which, like our approach, incorporates the HU ranges as an inductive bias.
Instead of utilizing binary masks, HU windowing creates multiple HU windows by clipping the HU values to each HU range and normalizing the values.
For a given tissue class $c$, let $L_{c}$ and $U_{c}$ denote the lower and upper bounds of the HU range, respectively. The transformed input representation is then defined as:
%
\begin{equation}
X_{c}(i,j) = \begin{cases}
    \frac{L_c}{|U_c - L_c|}, & \text{if } X_{0}(i,j) < L_{c},\\
    \frac{U_c}{|U_c - L_c|}, & \text{if } X_{0}(i,j) > U_{c},\\
    \frac{X_{0}(i,j)}{|U_c - L_c|}, & \text{otherwise}.
\end{cases}
\label{eq:pixel_clipping}
\end{equation}
%
% While HU windowing has been reported to improve segmentation performance, it remains relatively underexplored in practice. 
% To assess its efficacy compared to our proposed binary mask encoding, we conducted experiments using both approaches.
The results using HU windowing in \cref{tbl:hu-windowing} indicate that while our results are slightly superior, both inductive bias strategies yield comparable performance across different model sizes.
However, binary masks have a stronger inductive bias as HU values are \emph{reduced} to binary values, leading to fewer possible inputs.
Consequently, models trained with binary masks should learn simpler yet more robust patterns independent of the specific HU values.
\Cref{fig:hu-window-data} indicates that our approach works better in \emph{data-limited} circumstances, but further investigation is required to determine whether specific segmentation tasks or dataset characteristics favor one encoding method over the other.

\begin{table}[h]
    \caption{Results of HU windowing on our dataset.}
    \begin{center}
    \begin{tabular}{lr}
        \toprule
        \multicolumn{1}{c}{\textbf{Model}} 
        & \multicolumn{1}{c}{\textbf{Dice score}} \\
        \midrule
        U-Net-small, HU windows                    
        % & \formatpercentage{0.8804}\ \ $\pm$\formatpercentage{0.0048} \\               
        & \formatpercentage{0.9367}\ \ $\pm$\formatpercentage{0.004573} \\               
        \midrule
        U-Net-base,\ \ HU windows
        % & \formatpercentage{0.8840}\ \ $\pm$\formatpercentage{0.0066} \\               
        & \formatpercentage{0.9409}\ \ $\pm$\formatpercentage{0.003887} \\
        \midrule
        U-Net-large, HU windows                       
        % & \formatpercentage{0.8872}\ \ $\pm$\formatpercentage{0.0033} \\
        & \formatpercentage{0.9402}\ \ $\pm$\formatpercentage{0.002434} \\
        \bottomrule
    \end{tabular}
    \label{tbl:hu-windowing}
    \end{center}
\end{table}
\vspace{-2em}
\begin{figure}[h]
    \centering
    \includegraphics[width=1.0\linewidth]{figures/results/hu-window-data.png}
    \caption{\Cref{fig:robustness-data-size} with HU windowing experiments added. The vertical green line indicates that HU windowing surpasses the baseline performance with \SI{80}{\percent} of the dataset.}
    \label{fig:hu-window-data}
\end{figure}
